% main.tex — Dempsey Lab manuscript template (JASA-style article)
% GENERAL METHODS structure: Intro (w/ Contributions), Related Work,
% Method (Notation / Core method / Theory), Simulations, Application.
% For estimand-driven causal papers use main-causal.tex instead (Setup and
% Estimand / Identification / Estimation). Keep one, delete the other, and
% name the keeper main.tex. Comments explain what belongs in each section.
\documentclass[12pt]{article}

\usepackage{amsmath, amssymb, mathtools}
\usepackage{amsthm}
\usepackage[toc,page]{appendix}
\usepackage{graphicx}
\usepackage{placeins}
\usepackage{pgfplots}
\usepackage[shortlabels]{enumitem}
\usepackage{xcolor}
\usepackage{natbib}
\usepackage{float}
\usepackage[parfill]{parskip}
\usepackage{url}
\usepackage{hyperref}
\usepackage{booktabs}
\usepackage{bbm, bm}
\usepackage{amsfonts}

% Shared lab notation (operators, sequential-decision macros, theorem
% environments). Load BEFORE the newtx fonts. Do not redefine its macros.
\usepackage{lab-notation}

\usepackage{newtxtext}
\usepackage[subscriptcorrection]{newtxmath}

% DON'T change margins - should be 1 inch all around.
\addtolength{\oddsidemargin}{-.5in}%
\addtolength{\evensidemargin}{-.5in}%
\addtolength{\textwidth}{1in}%
\addtolength{\textheight}{-.3in}%
\addtolength{\topmargin}{-.8in}%

% NOTE: To produce blinded version, replace "0" with "1" below.
\newcommand{\blind}{0}

\def\spacingset#1{\renewcommand{\baselinestretch}{#1}\small\normalsize}
\spacingset{1}

\pgfplotsset{compat=1.18}

\begin{document}

\if1\blind
{
  \title{\bf [Project Title]}
  \author{Author 1\thanks{The authors gratefully acknowledge relevant funding sources in the unblinded version.}\\
    Department of Statistics}
  \maketitle
} \fi

\if0\blind
{
  \bigskip
  \bigskip
  \bigskip
  \begin{center}
    {\LARGE\bf [Project Title]}

    \medskip
    [DRI] and W. Dempsey

    Department of Biostatistics, University of Michigan, Ann Arbor, U.S.A.
  \end{center}
  \medskip
} \fi

\begin{abstract}
% One paragraph, written to be readable by someone who reads nothing else:
% the setting, the gap, the proposed estimand/method, what theory is
% established, and one sentence each on simulations and the application.
[Abstract placeholder.]
\end{abstract}

\noindent%
{\it Keywords:} [keyword 1]; [keyword 2]; [keyword 3].
\vfill

\newpage
\spacingset{1.45} % DON'T change the spacing!

\section{Introduction}
\label{sec:intro}
% The scientific setting, the problem, and why it matters. Motivate with a
% concrete example the reader can hold onto for the rest of the paper.
[Placeholder.]

\subsection{Contributions}
\label{sec:contributions}
% Numbered list of what this paper adds, each contribution one sentence,
% each verifiable (a theorem, an algorithm, an empirical finding). End
% with a one-paragraph roadmap of the paper.
[Placeholder.]

\section{Related Work}
\label{sec:related}
% Organize by theme, not by paper. Each paragraph: what that line of work
% achieves, and precisely what gap remains that Section~\ref{sec:method}
% fills. Be generous and accurate — reviewers come from these literatures.
[Placeholder.]

\section{Method}
\label{sec:method}

\subsection{Notation and problem formulation}
\label{sec:notation}
% Define the data structure, spaces, and target of inference BEFORE any
% method details. Use lab-notation.sty macros so symbols match other lab
% papers. State assumptions as numbered \begin{assumption} environments.
[Placeholder.]

\subsection{Core method}
\label{sec:core-method}
% The main construction: estimator, algorithm, or model. Lead with the
% idea in words, then the formal version. A boxed algorithm or displayed
% estimating equation the reader can implement is the goal.
[Placeholder.]

\subsection{Theory}
\label{sec:theory}
% Guarantees for the core method: consistency, rates, asymptotic
% normality, variance estimation. State results here; proofs go to the
% appendix. After each theorem, one paragraph on what it means in practice.
[Placeholder.]

\section{Simulation Study}
\label{sec:sims}
% One subsection per study, mirroring the simulations/ subdirectories in
% the repo — each answers a single named question (bias? coverage?
% efficiency gain? robustness to X?). State the data-generating process,
% the estimators compared, and the metric before showing tables.
[Placeholder.]

\section{Application}
\label{sec:application}
% The motivating data analysis. Describe the study, why the method fits,
% results, and a sensitivity analysis. Numbers here must be reproducible
% from the code repository.
[Placeholder.]

\section{Discussion}
\label{sec:discussion}
% Limitations, honest scope of the assumptions, and 2-3 concrete
% directions for future work.
[Placeholder.]

\bibliographystyle{agsm}
\bibliography{refs}

\begin{appendices}

\section{Proofs}
\label{app:proofs}
% One appendix section per main-text section that has results, in order.
[Placeholder.]

\section{Additional Simulation and Application Details}
\label{app:extra}
% Extra tables/figures, diagnostics, and anything reviewers may want but
% readers don't need.
[Placeholder.]

\section{Code and Data Availability}
\label{app:code}
% Link the GitHub repository and state where the data live and how access
% is governed (mirror data/README.md).
[Placeholder.]

\end{appendices}

\end{document}
